\bookchapter{Capítulo 53: Guardar datos y pedir cosas a la red}

\chapterintro{Una app que solo “mira” está muerta; la que guarda datos y se conecta a la red vive de verdad. En este capítulo conocerás el “trío” de la capa de datos de Android: DataStore (ajustes pequeños), Room (datos estructurados) y Retrofit (hablar con el servidor), más su “coordinador general”, las corrutinas. Al acabarlo podrás escribir apps de verdad que recuerdan los hábitos del usuario y refrescan contenido en vivo.}

\booksection{Cómo elegir dónde guardar: mira primero esta tabla}

\begin{center}
\begin{tabularx}{\linewidth}{|X|X|X|X|}
\hline
\textbf{Opción} & \textbf{Qué conviene guardar} & \textbf{Características} & \textbf{Dónde se usa en este libro} \tabularnewline
\hline
DataStore & Ajustes sueltos (tema, interruptores) & Pares clave-valor, sencillo y seguro & Los ajustes de la app de gastos del capítulo 54 \tabularnewline
Room & Datos estructurados (movimientos, usuarios) & Envoltura real de una base de datos SQLite, con consultas & Los movimientos del capítulo 54 \tabularnewline
Retrofit & (no se guarda: se trae de la red) Datos del servidor & El estándar de facto de las bibliotecas de peticiones HTTP & Los datos del tiempo del capítulo 55 \tabularnewline
\hline
\end{tabularx}
\end{center}

\begin{center}\small\color{gray}Reparto de tareas del trío de datos\end{center}

Grábate esta regla: “los ajustes pequeños van a DataStore, los datos en bloque van a Room y los datos de fuera viajan por Retrofit”. SharedPreferences, tan visto en tutoriales viejos, ya fue reemplazado por DataStore, y el SQL a mano sobre SQLite también cedió el puesto a Room: aprende directo lo nuevo y no tomes el camino viejo.

\booksection{DataStore: guardar ajustes pequeños}

Para recordar valores sueltos como “la moneda que eligió el usuario la última vez” o “si está activado el modo oscuro”, usa DataStore (versión Preferences):

\begin{lstlisting}[language=Kotlin]
val Context.dataStore by preferencesDataStore(name = "settings")

// escribir: guardar un valor
suspend fun saveTheme(isDark: Boolean) {
    context.dataStore.edit { prefs ->
        prefs[THEME_KEY] = isDark
    }
}

// leer: sacar el valor (Flow avisa solo de los cambios)
val themeFlow: Flow<Boolean> = context.dataStore.data
    .map { it[THEME_KEY] ?: false }
\end{lstlisting}

Fíjate en que devuelve un Flow, un flujo de datos que “fluye”. En cuanto el valor cambia, todos los puntos que lo observan reciben el nuevo valor sin que nadie les avise. Ese concepto (reactividad) es la piedra angular del Android moderno y lo llevarás a la práctica en el capítulo 54.

\booksection{Room: dar a los datos “casa propia”}

La app de gastos acumula cientos o miles de apuntes y necesita consultar y resumir por categoría y por mes: eso pide una base de datos. Room es la envoltura oficial de SQLite y deja “crear tablas y consultar” tan natural como escribir Kotlin. Son tres piezas:

\begin{lstlisting}[language=Kotlin]
// 1. Entity: la estructura de la tabla (una tabla = un data class)
@Entity(tableName = "records")
data class Record(
    @PrimaryKey(autoGenerate = true) val id: Long = 0,
    val amount: Double,
    val category: String,
    val note: String,
    val date: String
)

// 2. Dao: la interfaz de operaciones (alta, baja, cambio, consulta)
@Dao
interface RecordDao {
    @Query("SELECT * FROM records ORDER BY date DESC")
    fun getAll(): Flow<List<Record>>

    @Insert
    suspend fun insert(record: Record)

    @Delete
    suspend fun delete(record: Record)

    @Query("SELECT SUM(amount) FROM records WHERE category = :c")
    suspend fun sumByCategory(c: String): Double
}

// 3. Database: el ensamblaje final
@Database(entities = [Record::class], version = 1)
abstract class AppDatabase : RoomDatabase() {
    abstract fun recordDao(): RecordDao
}
\end{lstlisting}

Esta plantilla del “trío” puede mudarse tal cual a cualquiera de tus proyectos: basta con cambiar los nombres de los campos. Las consultas devuelven Flow y la interfaz se refresca al compás de los datos; no tiene precio.

\booksection{Retrofit: hablar con el servidor}

El tiempo, las noticias, los tipos de cambio... esos datos no viven en el móvil, sino en servidores ajenos. Retrofit deja pedir datos al servidor tan fácil como llamar a una función. Tres pasos:

\begin{lstlisting}[language=Kotlin]
// 1. definir la interfaz: describe las peticiones HTTP como funciones
interface WeatherApi {
    @GET("v1/weather")
    suspend fun getWeather(@Query("city") city: String): WeatherResponse
}

// 2. crear la instancia
val api = Retrofit.Builder()
    .baseUrl("https://api.example.com/")
    .addConverterFactory(GsonConverterFactory.create())
    .build()
    .create(WeatherApi::class.java)

// 3. llamar como si fuera una función (dentro de una corrutina)
val result = api.getWeather("Madrid")
\end{lstlisting}

\textbf{Aviso para esquivar problemas: dos leyes de hierro de la red}\par

(1) Las peticiones de red jamás en el hilo principal: congelarían la interfaz o provocarían un cierre (NetworkOnMainThreadException). Toda llamada de red debe ir dentro de una corrutina (viewModelScope.launch). (2) El HTTP en claro está prohibido por defecto: desde Android 9 se exige HTTPS. Si en pruebas necesitas una dirección local http://, tendrás que declarar usesCleartextTraffic a mano en el Manifest; pero un producto real debe ir todo por HTTPS.

\booksection{Parsear JSON: “traducir” los datos de la red}

La mayor parte de los datos que devuelve el servidor llega en JSON, un texto parecido a “diccionarios anidados”. Retrofit necesita un “intérprete” que lo convierta en objetos de Kotlin. Las dos opciones mayores de 2026:

\begin{center}
\begin{tabularx}{\linewidth}{|X|X|X|}
\hline
\textbf{Opción} & \textbf{Características} & \textbf{Escenario recomendado} \tabularnewline
\hline
Gson / Moshi & El veterano basado en reflexión, con coste de configuración cero & Proyectos pequeños de arranque rápido y con pocos campos \tabularnewline
kotlinx.serialization & La opción oficial en tiempo de compilación, más rápida y segura & Primera opción para proyectos nuevos (la corriente de 2026) \tabularnewline
\hline
\end{tabularx}
\end{center}

\begin{center}\small\color{gray}Las dos grandes opciones para parsear JSON\end{center}

\begin{lstlisting}[language=Kotlin]
// anota el data class con @Serializable
@Serializable
data class WeatherResponse(
    val weather: List<Weather>,
    val main: Main,
    val name: String
)

// Retrofit se le acopla el Converter correspondiente (ejemplo con el soporte oficial de Retrofit 2.11+):
// .addConverterFactory(json.asConverterFactory("application/json".toMediaType()))
\end{lstlisting}

Con cualquiera de las dos, la ley es la misma: mira primero el JSON real que llega y define los campos después. Muchos bugs de “fallo de parseo” vienen de nombres de campo mal escritos o de niveles de anidamiento mal adivinados; imprime el JSON en crudo y coteja: en diez segundos lo tienes.

\booksection{Corrutinas: el “idioma común” de la concurrencia en Android}

\booksection{Corrutinas: el “idioma común” de la concurrencia en Android}

Los suspend y launch que han ido saliendo arriba vienen de las corrutinas de Kotlin (Coroutine), la solución oficial de concurrencia de Android. Puedes entenderlo a lo burro así: hace que las operaciones “que obligan a esperar” (red, base de datos) no bloqueen la interfaz.

\begin{lstlisting}[language=Kotlin]
// arranca una corrutina dentro del ViewModel
viewModelScope.launch {
    val records = dao.getAll().first()  // suspend: se suspende y espera automáticamente
    // actualizar el estado de la interfaz…
}

// chuleta de palabras clave:
// suspend   —— "tengo que esperar, no bloquee a los demás"
// launch    —— "¡al trabajo!"
// viewModelScope —— el "taller" que sigue a la pantalla; se recoge solo al destruirse esta
\end{lstlisting}

Si partes de cero, con memorizar el uso basta; la teoría de “suspender y reanudar” que hay detrás puede esperar. El 90 \% de lo que escribirás a diario sigue esta plantilla.

\begin{tcbpractice}{Manos a la obra: crea una mini app de "frase del día"}
Busca una API de frases gratuita (algo tipo hitokoto), usa Retrofit para traer la “frase del día” y muéstrala en el centro de la pantalla, con un botón “Otra” que vuelva a pedir dentro de una corrutina. Después guarda las frases que te gusten en “favoritos” con Room. Este ejercicio pasa por DataStore/Room/Retrofit/corrutinas de arriba abajo; terminado, los capítulos 7 y 8 serán pan comido.

\end{tcbpractice}

\begin{tcbsummary}{Resumen del capítulo}
\begin{itemize}
  \item La regla: ajustes pequeños a DataStore, datos en bloque a Room, datos de fuera por Retrofit
  \item DataStore devuelve un Flow: si el valor cambia, los observadores lo reciben solos
  \item El trío de Room: @Entity (tabla) + @Dao (operaciones) + @Database (montaje final); la plantilla se recicla tal cual
  \item Retrofit en tres pasos: definir la interfaz crear la instancia llamar dentro de una corrutina
  \item Dos leyes de hierro: la red jamás en el hilo principal; producto real, siempre HTTPS
  \item Corrutina = esperar sin bloquear la interfaz; viewModelScope es el taller exclusivo de cada pantalla
\end{itemize}

\end{tcbsummary}
